\documentclass[noauthor,noinstructornotes]{ximera}

\title{Week 12}
\author{Math 1193 Design Team}

\begin{document}
\begin{abstract}
    Activities for Week 12.
\end{abstract}
\maketitle

\begin{problem}
\begin{enumerate}
    \item Draw the graph of \(g(x) = 3^x\). Be sure to label at least
    three points on the graph.
    \item Reflect your graph across the line \(y = x\). Call your new
    function \(f\).
    \begin{enumerate}
        \item What is the relationship between \(g\) and \(f\)?
        \item What points are on the graph of \(f\)?
        \item What is another name for \(f\)? There are two possible answers:
        see if you can come up with both!
    \end{enumerate}
    \item Is the asympote to the graph of \(f\) vertical or horizontal?
    What is its equation?
    \item What are the domain and range of \(f\)? If you get
    stuck, review the relationship between the domains and 
    ranges of inverse functions.
    
\end{enumerate}
\end{problem}

\begin{problem}
Aaliyah and Tom are having a friendly disagreement. Aaliyah thinks
\(\log(15) = \log(3) + \log(5)\) and Tom thinks
\(\log(15) = \log(3) \cdot \log(5)\).

\begin{itemize}
    \item Aaliyah's work is as follows:
    \begin{align*}
        \log(15) & = \log(3 \cdot 5) \\
        & = \log(10^{\log(3)} \cdot 10^{\log(5)}) \\
        & = \log(10^{\log(3) + \log(5)}) \\
        & = \log(3) + \log(5)
    \end{align*}
    \item Tom's work is as follows:
    \begin{align*}
        \log(15) & = \log(3 \cdot 5) \\
        & = \log(10^{\log(3)} \cdot 10^{\log(5)}) \\
        & = \log(10^{\log(3) \cdot \log(5)}) \\
        & = \log(3) \cdot \log(5)
    \end{align*}
\end{itemize}

Whose work is correct? Why? How could you check?

\end{problem}

\begin{problem}
When I was learning logarithms, there was no way to put \(\log_2(5)\)
into my calculator, since there was only a \(\log\) button and a
\(\ln\) button, without the option to directly input a base other
than 10 or \(e\). In this exercise, we'll learn how my generation
dealt with this problem. 

Make sure to write down your work at each step.

\begin{enumerate}
    \item Write the equation \(\log_2(5) = x\) in exponential form.
    \item Now take the \textbf{natural log} of both sides of your
    exponential equation.
    \item Next, we'll use a property of logarithms we'll learn soon:
    \(\ln(2^x) = x \ln(2)\). Use this to rewrite one side of your equation.
    \item Finally, solve for \(x\).
    \item How did I put \(\log_2(5)\) into my calculator? What about
    \(\log_7(19)\)?
\end{enumerate}
\end{problem}

\begin{problem}
\begin{enumerate}
    \item Sky was asked to simplify \(\ln(e^{-2})\) and said the answer was
    \(-2\). Are they correct? Why or why not? Write down your reasoning.
    \item Sky was asked to simplify \(e^{\ln(-2)}\) and said the answer was
    \(-2\). Are they correct? Why or why not? Write down your reasoning.
\end{enumerate}
\end{problem}

\end{document}
